\documentclass[11pt]{article}

\usepackage[a4paper,margin=1in]{geometry}
\usepackage{amsmath,amssymb,amsfonts,amsthm,mathrsfs}
\usepackage{enumitem}
\usepackage{booktabs}
\usepackage{xcolor}
\usepackage[hidelinks]{hyperref}

\emergencystretch=2em
\raggedbottom

\theoremstyle{plain}
\newtheorem{theorem}{Theorem}[section]
\newtheorem{proposition}[theorem]{Proposition}
\newtheorem{corollary}[theorem]{Corollary}
\newtheorem{lemma}[theorem]{Lemma}
\newtheorem{conjecture}[theorem]{Conjecture}

\theoremstyle{definition}
\newtheorem{definition}[theorem]{Definition}
\newtheorem{problem}[theorem]{Problem}
\newtheorem{example}[theorem]{Example}

\theoremstyle{remark}
\newtheorem{remark}[theorem]{Remark}

\newcommand{\Z}{\mathbb Z}
\newcommand{\Q}{\mathbb Q}
\newcommand{\R}{\mathcal R}
\newcommand{\G}{\mathcal G}
\newcommand{\A}{\mathcal A}
\newcommand{\F}{\mathcal F}
\newcommand{\I}{\mathcal I}
\newcommand{\Pset}{\mathcal P}
\newcommand{\Lsys}{\mathcal L}
\newcommand{\MSI}{\mathsf{MSI}}
\newcommand{\supp}{\operatorname{supp}}
\newcommand{\lcm}{\operatorname{lcm}}
\newcommand{\rad}{\operatorname{rad}}
\newcommand{\dist}{\operatorname{dist}}
\newcommand{\one}{\mathbf 1}

\title{The Goldbach Survivor-Gap Problem for Two-Cloud Primorial Masks\\
\large A Finite Jacobsthal-Type Reduction of Binary Goldbach\\[0.3em]
\normalsize Preprint version v0.4, June 2026 (10B computational validation)}
\author{Gabriel Dorel Dura\\Independent Researcher\\Braila, Romania\\\texttt{gabi.dura@gmail.com}}
\date{}

\begin{document}
\maketitle

\begin{abstract}
This paper isolates a single finite combinatorial problem naturally attached to
binary Goldbach.  For an even integer \(N\) and a sieve level \(z\), set
\[
Q(z)=\prod_{p\le z}p
\]
and define the Goldbach survivor mask
\[
\G_z(N)=\{x\bmod Q(z):(x,Q(z))=1,\ (N-x,Q(z))=1\}.
\]
Equivalently, for every prime \(p\le z\), the residue \(x\) avoids the two-cloud
forbidden set
\[
F_p(N)=\{0,\,N\bmod p\}.
\]
The Chinese remainder theorem gives the exact full-period count
\[
|\G_z(N)|=\prod_{p\le z}(p-c_p(N)),\qquad
c_p(N)=\begin{cases}1,&p\mid N,\\2,&p\nmid N.\end{cases}
\]
Thus the average spacing of survivors is of logarithmic-square scale, up to the
usual Goldbach singular factor.  However, the maximum gap cannot be of this
scale: complete sieving removes the consecutive interval \([2,z]\), so the
empty-run gap is at least \(z-1\).  The correct finite target is therefore not a
Cramer-type \((\log z)^2\) maximum-gap bound, but a Goldbach-scale subquadratic
bound.

Let \(g_z(N)\) denote the maximum length of a consecutive interval, modulo
\(Q(z)\), containing no element of \(\G_z(N)\).  The central problem formulated
here is to prove, at the Goldbach complete-sieve scale \(z\asymp \sqrt N\), that
\[
g_z(N)=o(N),
\]
or more concretely \(g_z(N)=o(z^2)\).  We prove the exact reduction: if this
subquadratic survivor-gap statement holds for all sufficiently large even
\(N\), then strong Goldbach follows for all sufficiently large even \(N\), with
only finite checking and explicit small-prime boundary channels left outside the
large-prime interior.  The proof is purely finite: a survivor
\(x\in(z,N-z)\) at level \(z\ge\sqrt N\) forces both \(x\) and \(N-x\) to be
prime.

The paper also records the survivor-polynomial formulation.  In
\(\Z[X]/(X^{Q(z)}-1)\), the mask polynomial
\[
S_{z,N}(X)=\sum_{x\in\G_z(N)}X^x
\]
factors into lifted local CRT masks.  For every interval \(I\), the localized
survivor count is the constant coefficient
\[
R_z(N;I)=[X^0]_{Q(z)}W_I(X)S_{z,N}(X^{-1}).
\]
Thus the gap problem is exactly a uniform coefficient-positivity problem for
interval polynomials.

The new material in this version formulates three concrete routes suggested by
this finite reformulation.  First, the autocorrelation of the two-cloud mask has
an exact CRT product formula, leading to a variance problem for interval counts.
Second, residual-prime distribution can be averaged over even parameters
\(N\) in a dyadic range, giving a realistic validation and partial-theorem target.
Third, survivor-polynomial truncations suggest a computational-theoretical search
for rigid finite minorants.  The paper also formulates the finite torus-minorant
problem: a pointwise lower-bound divisor weight, strong enough on the
complete-sieve interval, would imply binary Goldbach.  These are finite
combinatorial routes, not assumed analytic inputs.

The paper also records numerical evidence and a configurable computational
workflow.  In a high-range scan, every even \(N\le 10^{10}\) except the trivial
boundary-only case \(N=4\) had a complete-sieve interior representation; the
largest observed offset of the first interior coordinate above
\(\lfloor\sqrt N\rfloor\) was \(3464\).  These computations are validation
evidence for the formulation only.  They are not used as proof.

This paper does not prove Goldbach.  Its contribution is a precise finite
reduction and a sharply formulated candidate bridge: a subquadratic gap theorem
for Goldbach two-cloud primorial survivor masks.
\end{abstract}

\medskip
\noindent\textbf{Role in the series.}  This preprint is intended as the next
specialized paper in the MSI survivor-envelope program.  It takes the finite
survivor, windowed-envelope, boundary-channel, and survivor-polynomial tools of
the previous papers and specializes them to the binary Goldbach diagonal.  The
paper is a reduction and problem-formulation paper, not a proof of the binary
Goldbach conjecture.

\medskip
\noindent\textbf{Keywords.} Goldbach conjecture; finite sieve; survivor gap;
Jacobsthal problem; primorial modulus; Chinese remainder theorem; two-cloud
mask; survivor polynomial; coefficient positivity; additive number theory;
computational number theory.

\section{Introduction}

The binary Goldbach conjecture asks whether every even integer \(N>2\) is a sum
of two primes.  In the symmetric affine notation used in the previous papers of
this program, one sets
\[
 n=N-2,\qquad H=\{0\},
\]
and obtains the two-coordinate packet
\[
\Lsys_{n,\{0\}}(i)=\{i+1,\,n-i+1\}.
\]
Writing \(x=i+1\), this is simply
\[
 x,\qquad N-x.
\]
Thus the binary Goldbach problem is a localized diagonal nonemptiness problem:
one must find an \(x\) in the bounded interval \(1<x<N-1\) such that both
coordinates are prime.

The finite survivor, windowed-envelope, boundary-channel, and survivor-polynomial
machinery used here is inherited from the global survivor-envelope framework
of \cite{Dura2026SurvivorEnvelopes}.  The present paper specializes those tools
to the binary Goldbach two-cloud mask and isolates the remaining gap problem.

Classical treatments of Goldbach-type problems, sieve methods, and the general
distribution of primes provide important background and calibration
\cite{HardyLittlewood1923,HalberstamRichert1974,IwaniecKowalski2004,Koukoulopoulos2019}.
The present paper uses none of these analytic theories as proof input.  Its
statements below are finite congruence identities, conditional finite
reductions, or explicitly labelled finite problems.

The complete-sieve viewpoint makes the finite skeleton transparent.  If
\(z\ge\sqrt N\), then any composite integer \(m\le N\) has a prime divisor
\(p\le z\).  Hence any integer \(m\in(1,N]\) surviving all primes \(p\le z\) is
prime, except that prime values \(m\le z\) are removed by the complete sieve
itself and must be recorded as boundary channels.  Therefore binary Goldbach can
be separated into two finite pieces:
\[
\text{large-prime interior survivors}
\]
and
\[
\text{small-prime boundary channels.}
\]
This separation is exact and uses no analytic prime-distribution input.

The question is then whether the large-prime interior contains a survivor.  At
level \(z\), define
\[
Q(z)=\prod_{p\le z}p
\]
and the Goldbach survivor mask
\[
\G_z(N)=\{x\bmod Q(z):(x,Q(z))=1,\ (N-x,Q(z))=1\}.
\]
The mask is finite and periodic modulo \(Q(z)\).  It is described locally by the
forbidden residue set
\[
F_p(N)=\{0,N\bmod p\}\subset\Z/p\Z.
\]
For an even \(N\), local admissibility is automatic: modulo \(2\) the two
forbidden classes collide, while for odd \(p\) at most two classes are removed.

The central quantity in this paper is not the global size of \(\G_z(N)\), which
is an immediate CRT product, but the maximum empty interval in its support.  Let
\(g_z(N)\) denote the largest number of consecutive residue classes modulo
\(Q(z)\) containing no survivor.  If, at the complete-sieve scale \(z\asymp
\sqrt N\), one can prove
\[
g_z(N)=o(N),
\]
then the interval \((z,N-z)\) must contain a survivor for all sufficiently large
even \(N\).  This survivor gives a Goldbach representation with both primes
larger than \(z\).  Small-prime representations are handled separately by the
explicit channels \(x=q\) or \(N-x=q\), \(q\le z\).

This is the one-tool reduction proposed here:
\[
\boxed{\text{prove a subquadratic gap bound for Goldbach two-cloud survivor masks.}}
\]
The average spacing of \(\G_z(N)\) is of logarithmic-square size, but a maximum
gap of that size is impossible in the complete mask: the interval \([2,z]\) is
always removed by the condition \((x,Q(z))=1\).  The correct target is therefore
not \(g_z(N)\ll(\log z)^2\), but rather
\[
g_z(N)=o(z^2)
\]
at the Goldbach scale \(N\asymp z^2\).  This statement is much weaker than a
Cramer-type maximum-gap bound, but still strong enough to imply the binary
Goldbach conjecture outside a finite range.

The rest of the paper is organized as follows.  Section~\ref{sec:mask} defines
the two-cloud Goldbach mask and gives its exact CRT count.  Section~\ref{sec:gap}
defines the empty-run gap and explains the forced boundary lower bound.
Section~\ref{sec:goldbach_reduction} proves the finite reduction from a
subquadratic Goldbach-scale gap bound to binary Goldbach.  Section~\ref{sec:poly}
records the survivor-polynomial coefficient formulation.  Section~\ref{sec:jacobsthal}
places the problem beside the ordinary Jacobsthal problem.  Section~\ref{sec:crt_induction}
formulates the CRT lift induction and the missing distribution lemma.
Section~\ref{sec:minorants} formulates the finite torus-minorant problem.
Section~\ref{sec:routes} records three concrete routes suggested by the finite
reformulation: autocorrelation and variance, averaged residual-prime
distribution, and survivor-polynomial minorant search.  Section~\ref{sec:weak}
explains how ternary Goldbach serves as a known validation case for the
additive-ladder picture.  Section~\ref{sec:computational} lists the computational
tasks suggested by the reduction.

\section{The Goldbach two-cloud mask}
\label{sec:mask}

\begin{definition}[Primorial modulus]
For \(z\ge2\), let
\[
Q(z)=\prod_{p\le z}p.
\]
All primes in this product are rational primes.
\end{definition}

\begin{definition}[Goldbach survivor mask]
Let \(N\) be an even integer.  The level-\(z\) Goldbach survivor mask is
\[
\G_z(N)=\{x\bmod Q(z):(x,Q(z))=1,\ (N-x,Q(z))=1\}.
\]
Equivalently,
\[
x\in\G_z(N)
\]
if and only if, for every prime \(p\le z\),
\[
x\not\equiv 0\pmod p,
\qquad
x\not\equiv N\pmod p.
\]
\end{definition}

For each prime \(p\le z\), define the local forbidden cloud
\[
F_p(N)=\{0,N\bmod p\}\subset\Z/p\Z
\]
and the local survivor set
\[
G_p(N)=(\Z/p\Z)\setminus F_p(N).
\]
Let
\[
c_p(N)=|F_p(N)|.
\]
Then
\[
c_p(N)=
\begin{cases}
1,&p\mid N,\\
2,&p\nmid N.
\end{cases}
\]
In particular, for even \(N\), the prime \(2\) contributes one forbidden class,
not two.

\begin{proposition}[Exact CRT count]
\label{prop:crt_count}
For every even integer \(N\) and every \(z\ge2\),
\[
|\G_z(N)|=\prod_{p\le z}(p-c_p(N)).
\]
Equivalently,
\[
\frac{|\G_z(N)|}{Q(z)}
=
\prod_{\substack{p\le z\\p\mid N}}\left(1-\frac1p\right)
\prod_{\substack{p\le z\\p\nmid N}}\left(1-\frac2p\right).
\]
\end{proposition}

\begin{proof}
The condition defining \(\G_z(N)\) is a system of independent congruence
avoidance conditions modulo the primes \(p\le z\).  For each prime \(p\le z\),
there are \(p-c_p(N)\) admissible residue classes.  The Chinese remainder
theorem identifies a residue class modulo \(Q(z)\) with a tuple of residue
classes modulo the primes \(p\le z\).  Multiplying the local counts gives the
claim.
\end{proof}

\begin{remark}[Average spacing]
The product in Proposition~\ref{prop:crt_count} implies that the average
spacing between survivors in a complete period is approximately of
\((\log z)^2\) size, multiplied by the reciprocal of the usual Goldbach singular
factor.  This is only an average statement.  The maximum empty interval can be
much larger, and the complete sieve forces a boundary gap of length at least
\(z-1\), as shown below.
\end{remark}

At the Goldbach scale \(z\asymp\sqrt N\), the number \(N\) cannot be divisible by
many distinct primes \(p\le z\), because
\[
\prod_{\substack{p\le z\\p\mid N}}p\le N\asymp z^2.
\]
Thus the collision case \(p\mid N\), where the two forbidden classes collapse to
one, is sparse in the logarithmic sense
\[
\sum_{\substack{p\le z\\p\mid N}}\log p\le \log N.
\]
This scale restriction is important.  A uniform statement over arbitrary
residues \(N\bmod Q(z)\) would include degenerate cases such as
\(N\equiv0\bmod Q(z)\), where the mask reduces to the ordinary reduced residue
system.  The Goldbach problem only requires the diagonal regime
\(N\asymp z^2\).

\section{The survivor-gap function}
\label{sec:gap}

\begin{definition}[Empty-run gap]
Let \(S\subset\Z/Q\Z\) be nonempty.  Define \(g(S)\) to be the largest integer
\(L\ge0\) for which there exists an interval of \(L\) consecutive residue
classes modulo \(Q\) containing no element of \(S\).  For the Goldbach mask we
write
\[
g_z(N)=g(\G_z(N)).
\]
Thus every interval of length \(g_z(N)+1\) modulo \(Q(z)\) contains at least one
Goldbach survivor.
\end{definition}

This definition measures the longest empty run rather than the distance between
successive survivors.  The two conventions differ by one.  The present
empty-run convention is convenient because it translates directly into interval
nonemptiness.

\begin{proposition}[Forced boundary lower bound]
\label{prop:boundary_lower}
For every even \(N\) and every \(z\ge2\),
\[
g_z(N)\ge z-1.
\]
Consequently, no uniform estimate of the form
\[
g_z(N)\le C(\log z)^2
\]
can hold for the complete Goldbach survivor mask.
\end{proposition}

\begin{proof}
For every integer \(x\in\{2,3,\dots,z\}\), the number \(x\) has a prime divisor
\(p\le z\).  Hence \((x,Q(z))>1\), so \(x\notin\G_z(N)\).  Therefore the interval
\([2,z]\) contains no survivor and has length \(z-1\).  Since
\(z/(\log z)^2\to\infty\), a bound \(C(\log z)^2\) is impossible.
\end{proof}

\begin{remark}[Boundary channels]
The lower bound in Proposition~\ref{prop:boundary_lower} is not a defect in the
framework.  It is the finite trace of the small-prime boundary-channel
phenomenon.  Complete sieving at level \(z\) removes prime values \(q\le z\)
because such values are divisible by themselves.  Therefore a complete-sieve
Goldbach proof must combine large-prime interior survivors with explicit
small-prime boundary channels.
\end{remark}

The decisive scale is therefore not the average spacing.  At complete level
\(z\asymp\sqrt N\), the usable large-prime interior interval has length
approximately
\[
N-2z\asymp z^2.
\]
Thus the gap theorem needed for Goldbach is any estimate that is smaller than
this length, for example
\[
g_z(N)=o(z^2)
\]
or a more explicit estimate such as
\[
g_z(N)\ll z(\log z)^A.
\]
The second form would be much stronger than needed, but it is compatible with
the forced boundary gap \(g_z(N)\ge z-1\).

\section{Finite reduction to binary Goldbach}
\label{sec:goldbach_reduction}

\begin{definition}[Interior and boundary at level \(z\)]
Let \(N\) be even and \(z\ge2\).  The large-prime interior interval is
\[
I_z(N)=\{x\in\Z:z<x<N-z\}.
\]
The small-prime boundary channels are the possible representations
\[
N=q+(N-q),\qquad q\le z,
\]
where \(q\) is prime and \(N-q\) is prime.
\end{definition}

\begin{lemma}[Complete-sieve prime interior]
\label{lem:complete_interior}
Let \(N\) be even and let \(z\ge\sqrt N\).  If
\[
x\in I_z(N)
\]
and
\[
x\in\G_z(N),
\]
then both \(x\) and \(N-x\) are prime.
\end{lemma}

\begin{proof}
Since \(x\in I_z(N)\), both coordinates satisfy
\[
z<x<N-z,
\]
so in particular
\[
1<x<N,
\qquad
1<N-x<N.
\]
Suppose \(x\) were composite.  Then \(x\le N\), so \(x\) has a prime divisor
\(p\le\sqrt x\le\sqrt N\le z\).  This contradicts \((x,Q(z))=1\).  Hence \(x\)
is prime.  The same argument applies to \(N-x\), using
\((N-x,Q(z))=1\).
\end{proof}

\begin{proposition}[Boundary-interior equivalence]
\label{prop:boundary_interior}
Let \(N\) be an even integer and let \(z\ge\sqrt N\).  Then \(N\) has a Goldbach
representation if and only if at least one of the following holds:
\begin{enumerate}[label=(\roman*)]
\item there exists \(x\in I_z(N)\cap\G_z(N)\);
\item there exists a prime \(q\le z\) such that \(N-q\) is prime.
\end{enumerate}
\end{proposition}

\begin{proof}
If (i) holds, Lemma~\ref{lem:complete_interior} gives a representation by two
primes.  If (ii) holds, the representation is explicit.

Conversely, suppose \(N=a+b\) with \(a,b\) prime.  If one of \(a,b\) is at most
\(z\), then (ii) holds.  If both are greater than \(z\), then take \(x=a\).  Both
\(x\) and \(N-x=b\) are greater than \(z\), less than \(N-z\), and coprime to
\(Q(z)\).  Thus \(x\in I_z(N)\cap\G_z(N)\), so (i) holds.
\end{proof}

\begin{conjecture}[Goldbach-scale two-cloud gap conjecture]
\label{conj:subquadratic}
Let \(N\) range over even integers and set
\[
z_N=\lceil\sqrt N\rceil.
\]
Then
\[
g_{z_N}(N)=o(N)
\]
as \(N\to\infty\) through even integers.  Equivalently, since
\(z_N^2\asymp N\),
\[
g_{z_N}(N)=o(z_N^2).
\]
\end{conjecture}

\begin{theorem}[Conditional Goldbach implication]
\label{thm:conditional_goldbach}
If Conjecture~\ref{conj:subquadratic} holds, then every sufficiently large even
integer is a sum of two primes.  If, in addition, the remaining finite range is
checked computationally, then binary Goldbach follows for all even \(N>2\).
\end{theorem}

\begin{proof}
Let \(N\) be even and large, and set \(z=z_N=\lceil\sqrt N\rceil\).  The interior
interval
\[
I_z(N)=\{x\in\Z:z<x<N-z\}
\]
has length
\[
|I_z(N)|=N-2z-1+O(1)=N+O(\sqrt N).
\]
By Conjecture~\ref{conj:subquadratic}, for sufficiently large \(N\),
\[
g_z(N)<|I_z(N)|.
\]
Since every interval of length \(g_z(N)+1\) modulo \(Q(z)\) contains a survivor,
\(I_z(N)\) contains some \(x\in\G_z(N)\).  Lemma~\ref{lem:complete_interior}
then implies that both \(x\) and \(N-x\) are prime.  Thus \(N\) has a Goldbach
representation.  The finite remainder can be checked directly.
\end{proof}

\begin{remark}[No analytic input]
The implication in Theorem~\ref{thm:conditional_goldbach} is purely finite.  It
uses only divisibility, the complete-sieve threshold \(z\ge\sqrt N\), and an
empty-interval bound for \(\G_z(N)\).  No Hardy--Littlewood, circle-method,
Bombieri--Vinogradov, Elliott--Halberstam, or Bateman--Horn input is used.
\end{remark}

\section{Survivor polynomials and coefficient positivity}
\label{sec:poly}

The gap problem can be expressed exactly as a coefficient-positivity problem in
the cyclic group algebra.

Let
\[
\A_z=\Z[X]/(X^{Q(z)}-1).
\]
For each prime \(p\le z\), let \(e_p\in\Z/Q(z)\Z\) be the CRT idempotent
satisfying
\[
e_p\equiv1\pmod p,
\qquad
 e_p\equiv0\pmod q\quad(q\le z,
q\ne p).
\]
Define the local mask polynomial
\[
M_{p,N}(X)=\sum_{a\in G_p(N)}X^{a e_p}\in\A_z.
\]

\begin{proposition}[CRT mask factorization]
\label{prop:mask_factorization}
In \(\A_z\), the Goldbach survivor polynomial
\[
S_{z,N}(X)=\sum_{x\in\G_z(N)}X^x
\]
factors as
\[
S_{z,N}(X)=\prod_{p\le z}M_{p,N}(X).
\]
\end{proposition}

\begin{proof}
Expanding the product gives monomials with exponents
\[
\sum_{p\le z}a_p e_p\pmod{Q(z)},
\qquad a_p\in G_p(N).
\]
By the defining property of the idempotents, this exponent is the unique CRT
class whose reduction modulo \(p\) equals \(a_p\) for every \(p\le z\).  Thus the
set of exponents appearing in the product is exactly \(\G_z(N)\), each with
multiplicity one.
\end{proof}

For an interval \(I\subset\Z\), define the window polynomial
\[
W_I(X)=\sum_{x\in I}X^x\in\A_z.
\]
Let
\[
R_z(N;I)=|I\cap\G_z(N)|.
\]

\begin{proposition}[Interval coefficient identity]
\label{prop:coefficient_identity}
For every finite interval \(I\),
\[
R_z(N;I)=[X^0]_{Q(z)}W_I(X)S_{z,N}(X^{-1}),
\]
where \([X^0]_{Q(z)}\) denotes the coefficient of the identity class in
\(\Z[X]/(X^{Q(z)}-1)\).
\end{proposition}

\begin{proof}
Expanding gives
\[
W_I(X)S_{z,N}(X^{-1})=
\sum_{x\in I}\sum_{r\in\G_z(N)}X^{x-r}.
\]
The constant coefficient modulo \(Q(z)\) counts the pairs \((x,r)\) with
\[
x\equiv r\pmod{Q(z)}.
\]
This is exactly the number of integers \(x\in I\) whose residue class modulo
\(Q(z)\) lies in \(\G_z(N)\).
\end{proof}

\begin{corollary}[Gap as coefficient positivity]
\label{cor:gap_coeff}
The inequality
\[
g_z(N)<L
\]
is equivalent to the statement that for every interval \(I\) of length \(L\),
\[
[X^0]_{Q(z)}W_I(X)S_{z,N}(X^{-1})>0.
\]
\end{corollary}

Thus the Goldbach-scale two-cloud gap conjecture is a uniform localized
coefficient-positivity problem for the CRT-factored polynomial
\(S_{z,N}\).

\section{Relation with the Jacobsthal problem}
\label{sec:jacobsthal}

The ordinary Jacobsthal problem studies long intervals of integers all having a
nontrivial common factor with a fixed modulus.  For the primorial modulus
\(Q(z)\), this is the problem of bounding gaps in the reduced residue system
\[
(\Z/Q(z)\Z)^\times.
\]
It is a one-cloud problem: for each prime \(p\le z\), only the class
\[
x\equiv0\pmod p
\]
is forbidden.  Classical work of Iwaniec gives foundational bounds for this
problem \cite{Iwaniec1978}; modern large-prime-gap constructions, such as
\cite{FordGreenKonyaginTao2016,Maynard2016}, also use covering congruences and
sieving constructions that are conceptually adjacent to Jacobsthal-type
questions.

The Goldbach mask is different.  It is a two-cloud problem:
\[
x\not\equiv0\pmod p,
\qquad
x\not\equiv N\pmod p.
\]
The second cloud is not arbitrary; it is a translate of the first cloud by the
same global parameter \(N\).  This correlation is the structural feature that
one would have to exploit to beat a purely worst-case covering argument.

\begin{problem}[Goldbach two-cloud Jacobsthal problem]
\label{prob:two_cloud_jacobsthal}
At the Goldbach scale \(z=\lceil\sqrt N\rceil\), prove a subquadratic upper
bound
\[
g_z(N)=o(z^2)
\]
for the two-cloud survivor mask
\[
\G_z(N)=\{x\bmod Q(z):x\not\equiv0,N\pmod p\text{ for all }p\le z\}.
\]
\end{problem}

\begin{remark}[Why the scale restriction matters]
Problem~\ref{prob:two_cloud_jacobsthal} is not the same as a uniform bound over
all possible residues \(N\bmod Q(z)\).  Goldbach requires \(N\asymp z^2\).  This
prevents \(N\) from being divisible by an arbitrary set of small primes whose
product exceeds \(z^2\).  Therefore the true Goldbach problem is a restricted
family of two-cloud masks, not the full worst-case family of all masks modulo
\(Q(z)\).
\end{remark}

\section{CRT lift induction and the missing lemma}
\label{sec:crt_induction}

The natural induction adds primes one at a time.  Let
\[
Q_y=\prod_{p\le y}p
\]
and suppose \(p>y\) is the next prime.  Every residue class
\(a\bmod Q_y\) has \(p\) lifts modulo \(Q_y p\):
\[
a+tQ_y,
\qquad t=0,1,\dots,p-1.
\]
Since \(Q_y\) is invertible modulo \(p\), the congruence
\[
a+tQ_y\equiv0\pmod p
\]
removes exactly one value of \(t\), and the congruence
\[
a+tQ_y\equiv N\pmod p
\]
removes exactly one value of \(t\), unless the two values coincide.  Hence each
old survivor has at least \(p-2\) surviving lifts.

\begin{lemma}[Lift deletion bound]
\label{lem:lift_deletion}
Let \(a\bmod Q_y\) survive the Goldbach mask at level \(y\), and let \(p>y\) be
prime.  Among the \(p\) lifts
\[
a+tQ_y\bmod Q_y p,
\qquad 0\le t<p,
\]
at most two are removed by the new prime \(p\).  If \(p\mid N\), at most one is
removed.
\end{lemma}

\begin{proof}
The two new forbidden congruences are
\[
a+tQ_y\equiv0\pmod p,
\qquad
 a+tQ_y\equiv N\pmod p.
\]
Because \(Q_y\) is invertible modulo \(p\), each congruence has one solution in
\(t\bmod p\).  The two solutions coincide exactly when \(N\equiv0\pmod p\).
\end{proof}

Globally, Lemma~\ref{lem:lift_deletion} immediately recovers the product count.
Locally, it is not enough.  A short interval may see only a thin portion of the
lift grid, and the two deleted lift levels may erase all visible lifts of some
old residue classes.  The missing lemma must control this local distribution.

\begin{problem}[Residual-prime distribution lemma]
\label{prob:distribution_lemma}
Let \(y<z\), let \(I\) be an interval, and let
\[
\G_y(N)=\{x\bmod Q(y):(x,Q(y))=(N-x,Q(y))=1\}.
\]
For primes \(p\in(y,z]\), prove upper bounds of the form
\[
\#\{x\in I:x\bmod Q(y)\in\G_y(N),\ x\equiv0\text{ or }N\pmod p\}
\]
\[
\le
\left(\frac{2}{p}+\varepsilon_{y,p,I,N}\right)
\#\{x\in I:x\bmod Q(y)\in\G_y(N)\},
\]
with total error strong enough, after inclusion-exclusion or polynomial
coefficient control, to imply
\[
R_z(N;I)>0
\]
for every interval \(I\) of length \(o(z^2)\) at the Goldbach scale.
\end{problem}

A first-order union bound over the residual primes cannot be expected to close
the problem in full generality, because
\[
\sum_{y<p\le z}\frac{2}{p}
\]
can exceed \(1\).  Any successful proof must therefore use more structure.  The
survivor-polynomial identity provides the exact inclusion-exclusion algebra, and
the two-cloud relation \(\{0,N\}\) is the finite correlation that must replace a
purely arbitrary covering argument.


\section{Torus minorants and the finite error barrier}
\label{sec:minorants}

We now isolate the finite positivity problem that remains after the Goldbach
mask has been placed on the primorial torus.  The purpose of this section is to
separate completely proved finite identities from the additional minorant problem
that would be needed to close the Goldbach application.

Let \(d\mid Q(z)\) be squarefree.  Define
\[
\rho_d(N)=\prod_{p\mid d}c_p(N).
\]
This is the number of residue classes modulo \(d\) for which
\[
d\mid x(N-x).
\]

\begin{lemma}[Interval divisor-count lemma]
\label{lem:interval_divisor_count}
Let \(I\subset\Z\) be an interval.  For every squarefree \(d\mid Q(z)\),
\[
A_d(I,N)
:=
\#\{x\in I:d\mid x(N-x)\}
=
\frac{\rho_d(N)}{d}|I|+O(\rho_d(N)).
\]
The implied constant is absolute.
\end{lemma}

\begin{proof}
Modulo \(d\), the condition \(d\mid x(N-x)\) is the product of the local
conditions
\[
x\equiv0\pmod p
\quad\text{or}\quad
x\equiv N\pmod p
\qquad (p\mid d).
\]
By the Chinese remainder theorem, the number of such residue classes modulo
\(d\) is
\[
\rho_d(N)=\prod_{p\mid d}c_p(N).
\]
Each residue class modulo \(d\) occurs in the interval \(I\) as
\(|I|/d+O(1)\) integers.  Summing over the \(\rho_d(N)\) residue classes gives
the formula.
\end{proof}

\begin{proposition}[Exact inclusion--exclusion]
\label{prop:exact_ie}
For every interval \(I\),
\[
R_z(N;I)=\sum_{d\mid Q(z)}\mu(d)A_d(I,N).
\]
Consequently, the formal main term is
\[
|I|\prod_{p\le z}\left(1-\frac{c_p(N)}p\right).
\]
\end{proposition}

\begin{proof}
The condition defining \(R_z(N;I)\) is
\[
(x(N-x),Q(z))=1.
\]
Inclusion--exclusion over squarefree divisors of \(Q(z)\) gives
\[
\one_{(x(N-x),Q(z))=1}
=
\sum_{d\mid (x(N-x),Q(z))}\mu(d).
\]
Summing over \(x\in I\) gives the first identity.  Substituting the main term
from Lemma~\ref{lem:interval_divisor_count} gives
\[
|I|\sum_{d\mid Q(z)}\mu(d)\frac{\rho_d(N)}d
=
|I|\prod_{p\le z}\left(1-\frac{c_p(N)}p\right).
\]
\end{proof}

The difficulty is not the main term.  The difficulty is the accumulated error if
one uses the full inclusion--exclusion expansion directly.  A crude absolute
bound gives
\[
\sum_{d\mid Q(z)}\rho_d(N)
=
\prod_{p\le z}(1+c_p(N)),
\]
which is much larger than the expected main term at the complete Goldbach scale.
Thus the CRT identities alone do not prove positivity in the short interval
\[
I_z(N)=(z,N-z),
\qquad z=\lfloor\sqrt N\rfloor.
\]
This is the finite error barrier.

\subsection{A finite minorant problem}

The required strengthening may be formulated as a finite minorant problem.  We
seek coefficients \(\lambda_d\), supported on squarefree divisors \(d\mid Q(z)\),
such that
\[
\sum_{d\mid x(N-x)}\lambda_d
\le
\one_{(x(N-x),Q(z))=1}
\]
for every integer \(x\), and such that the corresponding interval sum is
positive:
\[
\sum_{x\in I_z(N)}\sum_{d\mid x(N-x)}\lambda_d>0.
\]
Since the minorant is pointwise bounded above by the survivor indicator,
positivity of this weighted sum implies
\[
R_z(N;I_z(N))>0.
\]

Using Lemma~\ref{lem:interval_divisor_count}, the weighted sum is
\[
\sum_{d\mid Q(z)}\lambda_d A_d(I_z(N),N)
=
|I_z(N)|\sum_{d\mid Q(z)}\lambda_d\frac{\rho_d(N)}d
+O\left(\sum_{d\mid Q(z)}|\lambda_d|\rho_d(N)\right).
\]
Thus a sufficient finite criterion is
\[
|I_z(N)|\sum_{d\mid Q(z)}\lambda_d\frac{\rho_d(N)}d
>
C\sum_{d\mid Q(z)}|\lambda_d|\rho_d(N),
\]
for an absolute constant \(C\), together with the pointwise minorant condition.

\begin{problem}[Rigid two-cloud minorant problem]
\label{prob:rigid_minorant}
For every even \(N>4\), with \(z=\lfloor\sqrt N\rfloor\), construct explicit
coefficients \(\lambda_d=\lambda_d(N,z)\), supported on squarefree \(d\mid Q(z)\),
such that
\[
\sum_{d\mid x(N-x)}\lambda_d
\le
\one_{(x(N-x),Q(z))=1}
\]
for every integer \(x\), and
\[
|I_z(N)|\sum_{d\mid Q(z)}\lambda_d\frac{\rho_d(N)}d
>
C\sum_{d\mid Q(z)}|\lambda_d|\rho_d(N).
\]
\end{problem}

This is a finite combinatorial problem.  It does not assume the
Hardy--Littlewood conjecture, the Bateman--Horn conjecture, the
Elliott--Halberstam conjecture, or any asymptotic prime-distribution estimate.

\begin{proposition}[Finite minorant implication]
\label{prop:minorant_implication}
If the rigid two-cloud minorant problem has a solution for an even integer
\(N>4\), then \(N\) has a Goldbach representation.
\end{proposition}

\begin{proof}
The minorant hypotheses imply
\[
0<\sum_{x\in I_z(N)}\sum_{d\mid x(N-x)}\lambda_d
\le
\sum_{x\in I_z(N)}\one_{(x(N-x),Q(z))=1}
=R_z(N;I_z(N)).
\]
Therefore there exists \(x\in I_z(N)\) such that
\[
(x,Q(z))=1,
\qquad
(N-x,Q(z))=1.
\]
Since \(x>z\) and \(N-x>z\), Lemma~\ref{lem:complete_interior} implies that
both \(x\) and \(N-x\) are prime.  Hence \(N=x+(N-x)\) is a Goldbach
representation.
\end{proof}

\subsection{The two-level tradeoff}

The existing full-period torus distribution can be localized only when the base
period is not too large compared with the interval.  Choose \(y<z\) and let
\[
Q(y)=\prod_{p\le y}p.
\]
If \(Q(y)\le |I_z(N)|\asymp z^2\), then the level-\(y\) torus gives meaningful
interval control.  Since \(Q(y)\approx e^y\), this requires roughly
\[
y\lesssim 2\log z.
\]
But then the residual primes \(y<p\le z\) have total first-order deletion mass
approximately
\[
\sum_{y<p\le z}\frac{2}{p}
\sim
2\log\frac{\log z}{\log y},
\]
which eventually exceeds \(1\).  Conversely, choosing \(y\) large enough to make
this residual sum small makes \(Q(y)\) much larger than the Goldbach interval,
so the current torus distribution theorem no longer localizes.

This tradeoff explains why the present theory does not yet prove the needed
gap bound.  The missing ingredient is not another full-period CRT count, but a
localized non-covering or minorant theorem exploiting the rigid relation
\[
F_p(N)=\{0,N\bmod p\}.
\]

\subsection{Relation with sieve weights}

The minorant formulation resembles Selberg-type lower-bound sieve constructions
and, more distantly, the weighted architectures used in modern prime-gap work
\cite{HalberstamRichert1974,Maynard2016}.  In the present paper this analogy is
only methodological.  No standard sieve theorem is invoked to prove
Problem~\ref{prob:rigid_minorant}.  The point is instead to identify the exact
finite inequality that a successful Goldbach-specific minorant would have to
satisfy.


\begin{remark}[Interpretation of the conditional implication]
The conditional implication above should not be read as making the Goldbach
problem disappear.  It relocates the difficulty into a finite complete-sieve gap
or coefficient-positivity statement.  At the scale
\[
z=\lfloor\sqrt N\rfloor,\qquad I_z(N)=(z,N-z),
\]
the assertion
\[
I_z(N)\cap\G_z(N)\ne\varnothing
\]
is precisely the large-prime-interior Goldbach assertion.  Thus a proof of a
subquadratic bound for the relevant maximum gap, or of a pointwise minorant that
forces this intersection, would be a proof of Goldbach in this finite language.
The value of the reformulation is that it makes the obstruction explicit and
suggests deterministic finite objects on which to work.
\end{remark}

\section{Three finite routes suggested by the reformulation}
\label{sec:routes}

The complete-sieve gap formulation is essentially equivalent in difficulty to
the large-prime-interior Goldbach problem.  Nevertheless, it exposes finite
structures that are less visible in the classical formulation.  We record three
routes that arise naturally from the survivor-polynomial and torus language.
They should be understood as attack routes on the remaining obstruction, not as
proofs of Goldbach.

\subsection{Autocorrelation and variance of the two-cloud mask}

Let
\[
f_{z,N}(x)=\one_{\G_z(N)}(x),\qquad
\mu_{z,N}=\frac{|\G_z(N)|}{Q(z)}.
\]
For an integer lag \(h\), define the full-period autocorrelation
\[
C_{z,N}(h)
=
\frac1{Q(z)}
\sum_{x\bmod Q(z)}f_{z,N}(x)f_{z,N}(x+h).
\]
The CRT gives an exact product formula.  For each prime \(p\le z\), let
\(\nu_p(h,N)\) be the number of distinct residue classes in
\[
0,\quad N,\quad -h,\quad N-h\pmod p.
\]
Then
\[
C_{z,N}(h)=\prod_{p\le z}\frac{p-\nu_p(h,N)}p.
\]
Indeed, the conditions \(x\in\G_z(N)\) and \(x+h\in\G_z(N)\) require \(x\) to
avoid the four classes listed above modulo each prime \(p\), with collisions
counted only once.

For a generic lag \(h\) and an odd prime \(p\nmid N\), the four classes are
distinct, so the local two-point factor is
\[
\frac{p-4}{p}.
\]
The square of the one-point local density is
\[
\left(\frac{p-2}{p}\right)^2.
\]
The generic correlation ratio is therefore
\[
\frac{(p-4)/p}{((p-2)/p)^2}
=
\frac{p(p-4)}{(p-2)^2}
=
1-\frac4{(p-2)^2}<1.
\]
Thus generic lags exhibit local anti-correlation.  Survivors do not behave like
independent random points at the local level; the two-cloud condition creates a
repulsive two-point factor for most lags.

For a cyclic interval of length \(M\), set
\[
A_M(a)=\sum_{0\le t<M} f_{z,N}(a+t).
\]
A standard expansion gives the exact variance identity
\[
\frac1{Q(z)}\sum_{a\bmod Q(z)}\bigl(A_M(a)-\mu_{z,N}M\bigr)^2
=
\sum_{|h|<M}(M-|h|)\bigl(C_{z,N}(h)-\mu_{z,N}^2\bigr),
\]
where the correlations are read modulo \(Q(z)\).  The main obstruction is the
set of exceptional lags for which congruences such as
\[
h\equiv 0,\ \pm N\pmod p
\]
hold for many small primes.  These lags create local collisions and can produce
positive correlations.  Controlling their total contribution strongly enough to
pass from average interval behavior to a maximum-gap bound is the first concrete
technical problem.

This route is the main direct attack on the complete-sieve gap problem: the
autocorrelation formula identifies the obstruction quantitatively, while the
variance identity gives a finite object that can be computed, estimated, and
compared with the observed complete-window gap data.

\subsection{Residual-prime distribution averaged over the parameter}

A second route is to relax the pointwise requirement in \(N\) and average over
even parameters in a dyadic range.  Let \(N\) vary over even integers in
\[
N\in[Z^2,2Z^2],
\]
with the complete level \(z\asymp Z\).  After sieving to an intermediate level
\(y<z\), one studies the base survivor set
\[
\G_y(N)=\{x:(x,Q(y))=1,\ (N-x,Q(y))=1\}
\]
inside the Goldbach window.  The residual clouds are
\[
x\equiv 0,N\pmod p,\qquad y<p\le z.
\]
For fixed \(N\), proving that these residual clouds do not cover all base
survivors is the hard pointwise problem.  When \(N\) varies, however, the
classes \(N\bmod p\) move through the residual primes, creating extra averaging
not present in the fixed-parameter problem.

A representative target is a mean-square residual distribution estimate of the
form
\[
\sum_{\substack{N\in[Z^2,2Z^2]\\ N\ \mathrm{even}}}
\left|
\#\{x\in I_z(N)\cap\G_y(N):x\equiv 0,N\pmod p\}
-\frac{c_p(N)}p\#(I_z(N)\cap\G_y(N))
\right|^2
\]
being small enough after summation over residual primes \(y<p\le z\), and then
over squarefree residual products.  Such estimates would measure whether the
residual clouds behave as predicted by the survivor-envelope model once the
translate parameter \(N\) is averaged.

This route is a validation and partial-theorem target.  Proving complete-sieve
positivity for a positive-density or density-one set of even \(N\) would not
prove Goldbach.  It would, however, demonstrate that the finite torus
reformulation produces distribution information at the Goldbach scale and would
isolate the exceptional pointwise obstruction more sharply.

\subsection{Rigid finite minorants from survivor-polynomial truncations}

The third route is a computational-theoretical search for a pointwise
certificate.  Instead of importing a generic lower-bound sieve, one starts from
the exact residual survivor-polynomial product.  After a base level \(y\), the
residual factor has the schematic form
\[
\prod_{y<p\le z}\bigl(I_p(X)-K_p(X)\bigr),
\]
where \(I_p(X)\) is the all-lift contribution and \(K_p(X)\) records the
forbidden local cloud at \(p\).  Expanding the product gives the exact residual
inclusion--exclusion.  The union-bound attrition estimate keeps only an
absolute first-order loss; the polynomial product retains all higher
intersections and alternating signs.

The proposed finite minorant search is as follows:
\begin{enumerate}[label=(\roman*),leftmargin=2em]
\item choose \(N,z,y\) and build the exact local residual masks;
\item truncate or smooth the residual product by conductor, support size,
residual degree, or squarefree level;
\item impose the pointwise minorant constraint
\[
T(x)\le \one_{(x(N-x),Q(z))=1};
\]
\item maximize the certified interval mass
\[
\sum_{x\in I_z(N)}T(x);
\]
\item search for stable coefficient patterns and then attempt to prove those
patterns symbolically.
\end{enumerate}

This differs from a standard Selberg lower-bound sieve in one essential respect:
the candidate weights are generated from the exact two-cloud survivor polynomial
rather than chosen only from generic divisor-count considerations.  The special
structure
\[
F_p(N)=\{0,N\bmod p\}
\]
is retained throughout.  The open question is whether this rigid translate
structure permits a finite pointwise minorant unavailable in a generic
dimension-two sieve.  This route is the most direct path toward a pointwise
Goldbach certificate, but it is also the route most exposed to the parity
barrier.

\subsection{Roles of the three routes}

The three routes play different roles in the program.  The autocorrelation
route is the main technical attack on the complete-sieve gap problem: it gives
an exact CRT formula for the two-point function and isolates the exceptional
lags that could support long gaps.  The averaged-over-\(N\) residual-prime route
is a validation and first partial-theorem target: it weakens the pointwise
problem while testing the predicted distribution at the Goldbach scale.  The
finite minorant route is a discovery mechanism for a possible pointwise proof:
survivor-polynomial truncations can be optimized computationally and then, if
stable patterns emerge, attacked symbolically.

Thus the proposed progression is
\[
\text{autocorrelation identifies the obstruction,}
\]
\[
\text{averaging over }N\text{ tests the predicted distribution,}
\]
\[
\text{and survivor-polynomial minorants search for a pointwise certificate.}
\]
None of these routes is presently a proof of Goldbach.  Together they specify
where the remaining finite work must occur.

\section{Weak Goldbach as a validation benchmark}
\label{sec:weak}

The ternary Goldbach theorem, proved by Helfgott \cite{Helfgott2013,Helfgott2015},
states that every odd integer greater than \(5\) is a sum of three primes.  In
the present paper this theorem is not used as an input to the binary problem.
It is used only as a validation benchmark for the additive-ladder viewpoint.

Binary Goldbach is a one-dimensional affine slice
\[
x_1+x_2=N
\]
inside a two-dimensional survivor space.  Ternary Goldbach is a two-dimensional
affine slice
\[
x_1+x_2+x_3=N
\]
inside a three-dimensional survivor space.  At a finite sieve level \(z\), the
ternary survivor count has the form
\[
R_z^{(3)}(N;I)=
\#\{(x_1,x_2,x_3)\in I^3:x_1+x_2+x_3=N,
\ (x_j,Q(z))=1\text{ for }j=1,2,3\}.
\]
At complete level, large-prime interior survivors give prime triples, while
coordinates \(x_j\le z\) are again explicit boundary channels.

The higher-dimensional slice has more room.  Therefore, in numerical and finite
survivor-envelope experiments, the ternary case should exhibit more stable
localized survivor mass than the binary diagonal.  This expectation agrees with
the historical and analytic fact that the ternary problem is now known, while the
binary problem remains open.  The validation role is modest but useful: it tests
whether the finite survivor machinery reflects the additive-dimensional hierarchy
\[
\text{unary primality}
\longrightarrow
\text{binary additive diagonals}
\longrightarrow
\text{higher-dimensional additive hyperplanes}.
\]


\section{Numerical validation and configurable computation}
\label{sec:numerics}

A numerical scan was run in the complete-sieve language described above.  The
first-interior part of the scan searched, for each even \(N\), for the first
prime \(x>\lfloor\sqrt N\rfloor\) such that \(N-x\) is also prime.  This is
exactly the large-prime interior condition from Lemma~\ref{lem:complete_interior}.

The high-range run used the dense backend and tested every even
\(N\le 10^{10}\).  Apart from the trivial boundary-only case \(N=4\), every
tested even integer had a complete-sieve interior representation.  The reported
classification was as follows:
\[
\begin{array}{rcl}
\text{even }N\text{ tested} &=& 4,999,999,999,\\
\text{interior positive} &=& 4,999,999,998,\\
\text{boundary-only} &=& 1,\\
\text{no representation found} &=& 0.
\end{array}
\]
The worst observed offset of the first interior coordinate above
\(\lfloor\sqrt N\rfloor\) was
\[
3464,
\]
occurring at
\[
N=9265570214,
\qquad
\lfloor\sqrt N\rfloor=96257,
\qquad
x=99721.
\]
Thus the first interior pair in that record offset case was
\[
99721+9265470493=9265570214.
\]
The largest first interior coordinate observed in the scan occurred at
\[
N=9997348832,
\qquad
\lfloor\sqrt N\rfloor=99986,
\qquad
x=102673,
\]
with offset \(2687\).  The worst ratio \(x/\lfloor\sqrt N\rfloor\) remains the
small case \(N=68\), where the first interior pair is \(31+37\).

A separate complete-window gap computation was performed exhaustively for every
even \(N\le250000\), and by sampling above that range.  In the exhaustive range,
the worst complete-window empty run was
\[
1439
\]
at
\[
N=137414,
\qquad z=370.
\]
In the high-range sampled data, the worst sampled complete-window empty run was
\[
6233
\]
at
\[
N=972350002,
\qquad z=31182.
\]
The ratio in this sampled case was approximately
\[
6.41\cdot 10^{-6}.
\]

These computations do not prove the gap conjecture.  They are included only as
validation evidence for the finite formulation: the complete-sieve interior is
not a rare numerical event, and the observed first-interior offset is very small
relative to \(N\) in the tested range.

A configurable C++ scanner, \texttt{Goldbach Lab v03}, was prepared to make the
numerical workflow reproducible and adjustable.  It supports a dense odd-only
prime-bitset backend for exhaustive scans within ordinary \(64\)-bit ranges, a
deterministic Miller--Rabin backend for sparse or very large tests, and separate
modes for first-interior scans, exact complete-window gap scans, and sampled gap
scans.  The intended output files are record cases, failure cases, sampled gap
records, and a summary report; full per-even-\(N\) output is deliberately avoided
for large limits.

The exhaustive first-interior scan to \(10^{10}\) was run with a command of the form
\[
\texttt{goldbach\_lab\_v03 --mode first-interior --maxN 10000000000 --backend dense}.
\]
The program is computational evidence infrastructure only.  It is not part of
the proof of any theorem in this paper.

\section{Computational tasks}
\label{sec:computational}

The reduction suggests a focused computational program for the MSI Survivor Lab.
For each even \(N\), set
\[
z=\lceil\sqrt N\rceil.
\]
Then compute or estimate:
\begin{enumerate}[label=(\arabic*)]
\item the local collision profile
\[
c_p(N)=1+\one_{p\nmid N},
\qquad p\le z;
\]
\item the global count
\[
|\G_z(N)|=\prod_{p\le z}(p-c_p(N));
\]
\item the first large-prime interior survivor in
\[
(z,N-z);
\]
\item the small-prime boundary channels \(q\le z\), \(N-q\) prime;
\item the largest observed empty run in restricted windows;
\item the layered certificate obtained by taking a base modulus \(Q(y)\le |I|\)
and then applying residual-prime attrition or polynomial inclusion-exclusion;
\item the ternary analogue as a validation comparison.
\end{enumerate}

A useful command-line workflow would classify each \(N\) into one of the
following categories:
\[
\begin{array}{c}
\text{interior survivor found},\\
\text{boundary representation found},\\
\text{layered certificate found},\\
\text{unexplained by current certificate}.
\end{array}
\]
The last category does not mean that Goldbach fails.  It means only that the
current finite certificate did not prove positivity for that instance.

\section{Conclusion}

The finite survivor framework turns binary Goldbach into a sharply formulated
support-spacing problem.  The relevant object is the two-cloud primorial mask
\[
\G_z(N)=\{x\bmod Q(z):(x,Q(z))=(N-x,Q(z))=1\}.
\]
At complete level \(z\ge\sqrt N\), any survivor in the interior interval
\((z,N-z)\) gives a Goldbach representation by two primes.  Therefore, a
Goldbach-scale subquadratic empty-run bound
\[
g_z(N)=o(z^2),
\qquad z\asymp\sqrt N,
\]
would imply binary Goldbach outside a finite range.

The paper does not prove this gap bound.  It identifies it as the central finite
problem and records the exact algebraic forms available for attacking it: CRT
mask factorization, interval coefficient positivity, complete-sieve boundary
separation, lift-by-prime induction, finite torus minorants, autocorrelation
variance identities, averaged residual-prime distribution problems, and
survivor-polynomial truncation searches.  The remaining task is to prove that
the correlated two-cloud forbidden structure \(\{0,N\bmod p\}\) cannot cover an
interval of Goldbach-scale length after the complete-sieve restriction is
imposed.

\begin{thebibliography}{99}

\bibitem{Dura2026Additive}
G. D. Dura,
\emph{Additive Encoding of Primality Constraints and Local Obstruction Structure in a Symmetric Affine System},
Zenodo preprint, 2026.
DOI: \href{https://doi.org/10.5281/zenodo.20259817}{10.5281/zenodo.20259817}.

\bibitem{Dura2026Packet}
G. D. Dura,
\emph{Exact Finite Packet Sieves and Survivor Convolutions for Symmetric Diagonal Affine Systems},
Zenodo preprint, 2026.
DOI: \href{https://doi.org/10.5281/zenodo.20573926}{10.5281/zenodo.20573926}.


\bibitem{Dura2026SurvivorEnvelopes}
G. D. Dura,
\emph{Global, Windowed, and Layered Survivor Envelopes: Primorial Survivor Spaces, Complete-Sieve Prime Interior, and Residual-Prime Attrition},
Zenodo preprint, 2026.
DOI: \href{https://doi.org/10.5281/zenodo.20617961}{10.5281/zenodo.20617961}.

\bibitem{FordGreenKonyaginTao2016}
K. Ford, B. Green, S. Konyagin, and T. Tao,
\emph{Large gaps between consecutive prime numbers},
Annals of Mathematics \textbf{183} (2016), 935--974.
DOI: \href{https://doi.org/10.4007/annals.2016.183.3.4}{10.4007/annals.2016.183.3.4}.

\bibitem{HalberstamRichert1974}
H. Halberstam and H.-E. Richert,
\emph{Sieve Methods},
London Mathematical Society Monographs, No. 4, Academic Press, 1974.

\bibitem{HardyLittlewood1923}
G. H. Hardy and J. E. Littlewood,
\emph{Some problems of `Partitio Numerorum'; III: On the expression of a number as a sum of primes},
Acta Mathematica \textbf{44} (1923), 1--70.

\bibitem{Helfgott2013}
H. A. Helfgott,
\emph{The ternary Goldbach conjecture is true},
arXiv:1312.7748, 2013.

\bibitem{Helfgott2015}
H. A. Helfgott,
\emph{The ternary Goldbach problem},
arXiv:1501.05438, 2015.

\bibitem{Iwaniec1978}
H. Iwaniec,
\emph{On the problem of Jacobsthal},
Demonstratio Mathematica \textbf{11} (1978), 225--231.

\bibitem{IwaniecKowalski2004}
H. Iwaniec and E. Kowalski,
\emph{Analytic Number Theory},
American Mathematical Society Colloquium Publications, Vol. 53, American Mathematical Society, 2004.

\bibitem{Koukoulopoulos2019}
D. Koukoulopoulos,
\emph{The Distribution of Prime Numbers},
Graduate Studies in Mathematics, Vol. 203, American Mathematical Society, 2019.

\bibitem{Maynard2016}
J. Maynard,
\emph{Large gaps between primes},
Annals of Mathematics \textbf{183} (2016), 915--933.
DOI: \href{https://doi.org/10.4007/annals.2016.183.3.3}{10.4007/annals.2016.183.3.3}.

\end{thebibliography}

\end{document}
